% Khai báo sử dụng AI và công cụ. Biên dịch bằng XeLaTeX: xelatex ai_declaration.tex (chạy hai lần).
\documentclass[11pt,a4paper]{article}

% Không dùng babel-vietnamese: gói này ép bảng mã T5 của pdfLaTeX, gây lỗi font với XeLaTeX.
\usepackage{fontspec}
\usepackage[margin=2.2cm]{geometry}
\usepackage{booktabs}
\usepackage{longtable}
\usepackage{array}
\usepackage{enumitem}
\usepackage[hidelinks]{hyperref}

% Nạp font theo tên tệp (có sẵn trong TeX Live/Overleaf, đủ dấu tiếng Việt), không phụ thuộc fontconfig.
\setmainfont{texgyretermes}[
  Extension = .otf,
  UprightFont = *-regular,
  BoldFont = *-bold,
  ItalicFont = *-italic,
  BoldItalicFont = *-bolditalic]
\renewcommand{\tablename}{Bảng}
\setlength{\emergencystretch}{3em}
\setlist{itemsep=2pt, topsep=4pt}
\newcolumntype{P}[1]{>{\raggedright\arraybackslash}p{#1}}

\title{Khai báo sử dụng AI và công cụ\\
\large Bài thực nghiệm giữa kỳ — Phân cụm khách hàng sàn thương mại điện tử}
\author{Mã sinh viên: 24022391}
\date{}

\begin{document}
\maketitle

\section{Công cụ AI}

\begin{itemize}
  \item \textbf{Trợ lý:} Claude Code (Anthropic), chạy dưới dạng tiện ích trong Visual Studio Code, có quyền đọc/ghi tệp trong thư mục dự án, chạy lệnh PowerShell/Python và thao tác git.
  \item \textbf{Mô hình:} Claude Sonnet 5.5 ở lượt đầu tiên (đọc và lưu các quy tắc). Sau đó sinh viên chuyển sang Claude Opus 5.5 bằng lệnh \texttt{/model opus}, và mô hình này được dùng cho toàn bộ phần thực hiện câu a–e, kiểm chứng và báo cáo.
  \item \textbf{Bộ nhớ của trợ lý:} các quy tắc chung của đề bài (vai trò, thông số, công nghệ, cấu trúc đầu ra, phong cách, kiểm chứng, báo cáo, quy trình) được lưu vào bộ nhớ dự án của Claude Code để áp dụng xuyên suốt.
\end{itemize}

\section{Phạm vi hỗ trợ}

\begin{longtable}{P{0.3\textwidth}P{0.64\textwidth}}
\toprule
\textbf{Bên thực hiện} & \textbf{Nội dung} \\
\midrule
\endhead
Sinh viên &
Cung cấp đề bài và các yêu cầu theo từng mục (1–13); đặt các ràng buộc (không ngụy tạo số liệu, không dùng nhãn giả lập khi huấn luyện, commit theo từng phần); chọn cách làm khi được hỏi (push trực tiếp lên nhánh \texttt{main}); yêu cầu kiểm chứng, báo cáo LaTeX, khai báo và rà soát mã. \\
\addlinespace
Trợ lý AI &
Đề xuất thiết kế dữ liệu, cơ chế lỗi và quy trình xử lý; viết toàn bộ mã Python trong notebook; chạy notebook, đọc kết quả và sửa lỗi phát sinh; viết nhận xét sau khi có kết quả chạy; viết báo cáo LaTeX và tệp khai báo này; tạo commit và push lên GitHub theo yêu cầu. \\
\bottomrule
\end{longtable}

Các nguyên tắc đã áp dụng trong quá trình làm:
\begin{itemize}
  \item Số liệu trong notebook và báo cáo đều lấy trực tiếp từ kết quả chạy. Báo cáo dùng macro LaTeX do notebook sinh ra, và notebook có kiểm tra tự động đối chiếu macro với các tệp kết quả.
  \item Các kết quả trái kỳ vọng được giữ nguyên và giải thích, ví dụ: điền khuyết theo nhóm không tốt hơn trung vị toàn cục; Davies–Bouldin xấu hơn sau tiền xử lý ở $k=4$.
  \item Dữ liệu gốc và metadata chỉ dùng để đánh giá sau khi xử lý xong, không dùng để chọn phương pháp hay tham số.
\end{itemize}

\section{Lịch sử hội thoại chính}

Bảng dưới tóm tắt các lượt trao đổi chính giữa sinh viên và trợ lý theo thứ tự thời gian (ngày 13/11/2026). Mã commit tham chiếu kho \texttt{github.com/Longkyoladn/Datamining-}.

\begin{longtable}{P{0.04\textwidth}P{0.38\textwidth}P{0.42\textwidth}P{0.08\textwidth}}
\toprule
\textbf{\#} & \textbf{Yêu cầu của sinh viên} & \textbf{Trợ lý thực hiện} & \textbf{Commit} \\
\midrule
\endhead
1 & Gửi các quy tắc chung (mục 1–3, 9–13), yêu cầu đọc và lưu, chưa tác động tới dự án. &
Lưu quy tắc vào bộ nhớ; nhận thấy mục 4–8 chưa có và hỏi lại. & — \\
\addlinespace
2 & Xác nhận mục 4–8 sẽ được gửi theo từng phần; chuyển mô hình sang Opus; gửi mục 4 (câu a). &
Kiểm tra môi trường Python; thiết kế 4 nhóm tiềm ẩn và cơ chế lỗi; viết, chạy notebook câu a; lưu dữ liệu original, corrupted, metadata; kiểm tra bằng \texttt{assert} và kiểm tra tái lập. & \texttt{205ade0} \\
\addlinespace
3 & Gửi mục 5 (câu b). &
EDA trên dữ liệu corrupted; quy tắc nghiệp vụ, IQR, đánh giá khả năng phát hiện; phân tích 4 vấn đề ảnh hưởng K-Means. & \texttt{6213fab} \\
\addlinespace
4 & Yêu cầu commit lên git theo từng phần. &
Hỏi nhánh đích; sinh viên chọn push trực tiếp lên \texttt{main}. Tách commit câu a (notebook chỉ gồm câu a) và câu b. & (2 commit trên) \\
\addlinespace
5 & Hỏi dữ liệu ban đầu, dữ liệu có nhiễu/ngoại lai và dữ liệu sau xử lý đã có chưa, được xử lý thế nào. &
Giải thích: dữ liệu original/corrupted đã có; dữ liệu preprocessed/prepared thuộc câu c–d, chưa tạo. & — \\
\addlinespace
6 & Gửi mục 6 (câu c). &
Giá trị không hợp lệ $\rightarrow$ NaN; điền khuyết; IQR sau sửa; không xóa bản ghi; bảng so sánh trước/sau; đánh giá sai số điền khuyết sau xử lý. & \texttt{1470fe9} \\
\addlinespace
7 & Gửi mục 7 (câu d). &
Bỏ \texttt{total\_spend\_6m}; log1p; RobustScaler; OneHotEncoder; PCA để trực quan hóa; lưu ma trận đặc trưng. & \texttt{1230455} \\
\addlinespace
8 & Gọi \texttt{/code-review} rồi ngắt; gửi mục 8 (câu e). &
Tách bộ biến đổi thành fit/transform dùng chung cho hai nhánh; khảo sát $k$; K-Means hai nhánh và tham chiếu sạch; ARI sau phân cụm; hồ sơ cụm, đặt tên, chiến lược tiếp thị. & \texttt{0760e36} \\
\addlinespace
9 & Yêu cầu kiểm tra pipeline theo mục 11, viết báo cáo LaTeX, tệp khai báo riêng, sau đó chạy \texttt{/code-review} và rà soát. &
Thêm 12 hàm kiểm chứng vào notebook; sinh macro và bảng LaTeX từ kết quả; sửa một giá trị làm tròn sai trong nhận xét câu e (1.897 $\rightarrow$ 1.896); viết \texttt{report.tex}, tệp này và \texttt{README.md}; chạy hai vòng rà soát mã và sửa các phát hiện (Mục~\ref{sec:review}). & \texttt{5b244f7}, \texttt{b80e3f6}, \texttt{9495d2b} \\
\addlinespace
10 & Báo lỗi định dạng khi biên dịch bằng XeLaTeX. &
Không biên dịch thử được (máy không có LaTeX, lệnh tải trình biên dịch bị chặn). Bỏ babel-vietnamese, nạp font TeX Gyre Termes theo tên tệp, sửa đường dẫn cho Overleaf. & \texttt{11fe3f5} \\
\addlinespace
11 & Yêu cầu trợ lý đóng vai giáo viên nghiêm khắc chấm điểm và chỉ ra hạn chế của báo cáo. &
Ước lượng khoảng 7.5/10; liệt kê lỗi nặng (thiếu report.md, một câu kết luận sai về điền khuyết, thiếu bảng so sánh trước/sau, hồ sơ cụm thiếu chi tiêu trung bình, thiếu chiến lược cho $k$ chính), các hạn chế phương pháp và trình bày. & — \\
\addlinespace
12 & Yêu cầu xuất nhật ký hội thoại dạng Markdown. &
Tạo \texttt{reports/conversation\_log.md} (tệp này sau đó không còn trong thư mục dự án và không được commit). & — \\
\addlinespace
13 & Yêu cầu sửa phần tiền xử lý và các vấn đề báo cáo liên quan. &
Sửa câu ``không kém hơn median imputation'' thành ``kém hơn một chút'' và thêm kiểm tra tự động; thêm ba bảng so sánh trước/sau tiền xử lý vào báo cáo; nêu rõ hạn chế của đánh giá điền khuyết và các giả định nghiệp vụ. & (commit cuối) \\
\bottomrule
\end{longtable}

\section{Công cụ và phiên bổ sung}

\begin{longtable}{P{0.3\textwidth}P{0.64\textwidth}}
\toprule
\textbf{Công cụ / phiên} & \textbf{Cách sử dụng} \\
\midrule
\endhead
Python 3.13.9 & numpy 2.5.3, pandas 2.3.3, matplotlib 3.10.7, seaborn 0.13.2, scikit-learn 1.7.2. scipy 1.16.3 có trong môi trường nhưng không được dùng. \\
\addlinespace
Jupyter kernel & ipykernel 6.30.1, jupyter\_client 8.6.3. Môi trường không có nbconvert/nbclient nên trợ lý dùng một script nhỏ dựa trên \texttt{jupyter\_client} để chạy notebook từ đầu đến cuối trong kernel mới và ghi lại output. \\
\addlinespace
Script dựng notebook & Trợ lý soạn nội dung các cell trong một script Python tạm (\texttt{nbformat}) ngoài thư mục dự án, rồi sinh ra \texttt{customer\_segmentation.ipynb}. Script này không nằm trong kho mã; notebook là sản phẩm chính và chạy độc lập. \\
\addlinespace
Kiểm tra tái lập ngoài notebook & Sau mỗi phần, notebook được chạy lại và mã băm (hash) của các tệp CSV được so sánh với lần chạy trước; các tệp trùng khớp. \\
\addlinespace
git / GitHub & Mỗi câu một commit, push lên nhánh \texttt{main} theo lựa chọn của sinh viên. \\
\addlinespace
Bộ nhớ Claude Code & Lưu các quy tắc chung của đề bài và phong cách viết để áp dụng cho mọi phần. \\
\addlinespace
\texttt{/code-review} (lần 1) & Sinh viên gọi trước câu e rồi ngắt; không có kết quả rà soát. \\
\addlinespace
\texttt{/code-review} (lần 2) & Chạy sau khi hoàn thành báo cáo; kết quả ở Mục~\ref{sec:review}. \\
\addlinespace
LaTeX & Máy thực hiện không có trình biên dịch LaTeX. Các tệp \texttt{.tex} được kiểm tra cú pháp bằng script (cân bằng ngoặc, cặp \texttt{begin}/\texttt{end}, macro được định nghĩa) nhưng chưa được biên dịch thành PDF trong phiên này. \\
\bottomrule
\end{longtable}

\section{Kết quả rà soát mã}
\label{sec:review}

Sinh viên yêu cầu chạy \texttt{/code-review} với tham số \texttt{-mode 5.6 sol - maxround 2}. Công cụ rà soát của Claude Code không định nghĩa các tham số \texttt{-mode} và \texttt{sol}, nên rà soát chạy ở mức mặc định (đọc mã, không chạy thử trên dữ liệu thay đổi). Giới hạn hai vòng được áp dụng: rà soát $\rightarrow$ sửa $\rightarrow$ rà soát lại $\rightarrow$ sửa, không có vòng thứ ba.

\subsection*{Vòng 1 — rà soát commit \texttt{5b244f7} (10 phát hiện, sửa ở \texttt{b80e3f6})}
\begin{itemize}
  \item Bảng kiểm chứng luôn ghi PASS; một kiểm tra lỗi làm dừng cell và để lại bảng PASS cũ. \emph{Sửa:} ghi nhận FAIL, luôn ghi bảng trước khi báo lỗi.
  \item Notebook lỗi khi $k$ chính trùng $k=4$. \emph{Sửa:} luôn chạy thí nghiệm $k=4$.
  \item Kích thước cụm nhỏ nhất hiển thị dạng số thực trong bảng khảo sát $k$. \emph{Sửa:} ép kiểu số nguyên.
  \item Báo cáo còn số liệu và nhận định gõ tay. \emph{Sửa:} thay bằng macro; thêm kiểm tra tự động cho các nhận định định tính.
  \item Tiêu đề bảng phát hiện ghi cứng số lượng và gán theo vị trí cột. \emph{Sửa:} đổi tên theo tên cột, số lượng lấy từ tham số.
  \item Kiểm tra 12 chỉ đối chiếu một phần macro và gần như vòng tròn. \emph{Sửa:} lưu nhãn cụm, tính lại các chỉ số từ tệp đã lưu, kiểm tra mọi macro dùng trong báo cáo đã được định nghĩa.
  \item Kiểm tra tái lập bỏ qua nhánh tham chiếu sạch; kiểm tra lưu tệp thiếu một số tệp; \texttt{latex\_escape} chưa đầy đủ; một số phép tính lặp và chuỗi ghi cứng. \emph{Sửa:} bổ sung tương ứng.
\end{itemize}

\subsection*{Vòng 2 — rà soát commit \texttt{b80e3f6} (10 phát hiện, sửa ở commit cuối)}
\begin{itemize}
  \item Kiểm tra nhận định định tính dùng \texttt{assert} trước khi ghi tệp, nên vẫn có thể để lại bảng PASS cũ. \emph{Sửa:} chỉ ghi nhận, kết quả FAIL được ghi ở kiểm tra 12 sau khi mọi tệp đã được ghi.
  \item Nhận định được kiểm tra trên dữ liệu trong bộ nhớ (luôn đúng). \emph{Sửa:} các nhận định về phân cụm được tính lại từ tệp đã lưu.
  \item Macro Silhouette/Davies–Bouldin không còn được đối chiếu. \emph{Sửa:} đối chiếu toàn bộ với \texttt{clustering\_metrics.csv} và tính lại Davies–Bouldin của nhánh B từ tệp.
  \item Trường hợp $k$ chính bằng 4 vẫn làm nhận định so sánh thất bại. \emph{Xử lý:} khi đó văn bản báo cáo không còn đúng, và kiểm tra 12 ghi FAIL (không còn dừng notebook trước khi ghi tệp).
  \item Báo cáo còn ghi cứng ``$k=2$'', tên bản ghi tạo cụm 1 phần tử, và mô tả Elbow chưa khớp với điều kiện kiểm tra. \emph{Sửa:} dùng macro \texttt{KMain}/\texttt{KDesign}; thêm nhận định $k$ chính bằng 2 và nhận định cụm 1 phần tử là bản ghi có tỷ lệ hoàn trả $>1$; viết lại câu về Elbow cho khớp điều kiện.
  \item Biểu thức tìm macro bắt nhầm lệnh LaTeX chuẩn viết hoa và phần chú thích. \emph{Sửa:} bỏ chú thích, loại các lệnh chuẩn.
  \item Nhãn cụm của tham chiếu sạch được ghép theo vị trí dòng mà không kiểm tra thứ tự. \emph{Sửa:} \texttt{assert} thứ tự \texttt{customer\_id} trùng nhau.
  \item Độ lệch chuẩn nhiễu dạng số bị trộn vào dict mô tả. \emph{Sửa:} trả về thành giá trị riêng.
\end{itemize}

\subsection*{Kiểm chứng sau rà soát}
Sau vòng 2, notebook được chạy lại hai lần: 12/12 kiểm tra đạt, 16/16 nhận định định tính đúng, và 19 tệp dữ liệu, kết quả, LaTeX sinh ra trùng khớp giữa hai lần chạy. Đường ghi FAIL được thử bằng cách tạm thêm một macro không tồn tại vào \texttt{report.tex}: kiểm tra 12 ghi FAIL kèm lý do, và cả tệp CSV lẫn bảng LaTeX đều được cập nhật. Sau khi khôi phục \texttt{report.tex}, chạy lại cho 12/12 đạt. Các trường hợp dữ liệu khác (seed khác, tham số khác) chưa được chạy thử.

\end{document}
